% =====================================================================================================
%  Estilo del manual de TurbuLENS: colores de la aplicación, recuadros, botones, capturas con marcas
%  numeradas y fragmentos de código con sus números de línea reales.
% =====================================================================================================
\usepackage{fontspec}
\usepackage{polyglossia}
\setdefaultlanguage{spanish}
\setotherlanguage{english}
\IfFontExistsTF{Inter}
  {\setmainfont{Inter}[Scale=0.92]\setsansfont{Inter}[Scale=0.92]}
  {\IfFontExistsTF{Roboto}{\setmainfont{Roboto}\setsansfont{Roboto}}{}}
\IfFontExistsTF{DejaVu Sans Mono}{\setmonofont{DejaVu Sans Mono}[Scale=0.8]}{}
\usepackage{amsmath,amssymb}
\usepackage{newunicodechar}
% Símbolos de los botones de la aplicación que la fuente del texto no tiene
\IfFontExistsTF{DejaVu Sans}{\newfontfamily\simbolos{DejaVu Sans}[Scale=0.85]}{\let\simbolos\relax}
\newunicodechar{◎}{{\simbolos ◎}}
\newunicodechar{⋯}{{\simbolos ⋯}}
\newunicodechar{★}{{\simbolos ★}}
\newunicodechar{■}{{\simbolos ■}}
\newunicodechar{▶}{{\simbolos ▶}}
\usepackage[a4paper,top=2.5cm,bottom=2.4cm,left=2.5cm,right=2.5cm,headheight=16pt]{geometry}
\usepackage{graphicx,xcolor,booktabs,tabularx,array,float,multicol,caption,fancyhdr,titlesec,ragged2e}
\usepackage[shortlabels]{enumitem}
\usepackage[most,listings]{tcolorbox}
\usepackage{tikz}
\usetikzlibrary{positioning,arrows.meta,calc,decorations.pathmorphing,shapes.misc,fit,backgrounds}
\usepackage{pgfplots}
\pgfplotsset{compat=1.16}
\usepackage{listings}
\usepackage[hyphens]{url}
\usepackage[hidelinks,bookmarksnumbered]{hyperref}

% ---------------------------------------------------------------------------------------------- colores
\definecolor{acento}{HTML}{3E7C72}      % verde azulado del tema claro
\definecolor{acentoosc}{HTML}{2A5A52}
\definecolor{naranja}{HTML}{F28C28}     % naranja del tema oscuro
\definecolor{cobre}{HTML}{B4652A}
\definecolor{suave}{HTML}{E8F1EF}
\definecolor{crema}{HTML}{FFF6EC}
\definecolor{gris}{HTML}{5B6661}
\definecolor{grisclaro}{HTML}{EEF0EF}
\definecolor{tinta}{HTML}{1F2623}
\definecolor{codigofondo}{HTML}{F6F8F7}
\definecolor{grisflecha}{HTML}{9AA39F}
\definecolor{grisnum}{HTML}{8E9893}
\hypersetup{colorlinks=true,linkcolor=acentoosc,urlcolor=acento,citecolor=acento,
  pdftitle={Manual de TurbuLENS},pdfauthor={Diego Fernando Cagua Lozada, Fredy Esteban Cruz Ospina, Saray Juliana Pinto Sanabria, José Alejandro López Hernández}}
\color{tinta}

% ---------------------------------------------------------------------------------------------- párrafos
\setlength{\parindent}{0pt}
\setlength{\emergencystretch}{2.5em}
\setlength{\parskip}{0.55em}
\RaggedRight
\setlist{itemsep=0.2em,topsep=0.35em,parsep=0.1em}
\newcolumntype{Y}{>{\RaggedRight\arraybackslash}X}
\renewcommand{\arraystretch}{1.22}
\captionsetup{font=small,labelfont={bf,color=acento},labelsep=period,justification=RaggedRight,singlelinecheck=false}
\setcounter{tocdepth}{1}
\addto\captionsspanish{\renewcommand{\contentsname}{Contenido}\renewcommand{\tablename}{Tabla}}
\setcounter{secnumdepth}{2}

% ---------------------------------------------------------------------------------------------- títulos
\titleformat{\chapter}[display]
  {\normalfont\bfseries\color{tinta}}
  {\begin{tikzpicture}[baseline]
     \node[fill=acento,text=white,rounded corners=3pt,inner sep=6pt,font=\Large\bfseries] {\thechapter};
   \end{tikzpicture}}
  {0.6em}{\fontsize{26}{30}\selectfont}
  [\vspace{0.2em}{\color{acento}\titlerule[1.2pt]}]
\titlespacing*{\chapter}{0pt}{-12pt}{18pt}
\titleformat{\section}{\Large\bfseries\color{acentoosc}}{\thesection}{0.7em}{}
\titleformat{\subsection}{\large\bfseries\color{tinta}}{\thesubsection}{0.7em}{}
\titleformat{\subsubsection}{\normalsize\bfseries\color{gris}}{}{0em}{}
\titlespacing*{\section}{0pt}{1.4em}{0.5em}
\titlespacing*{\subsection}{0pt}{1.1em}{0.35em}
\titlespacing*{\subsubsection}{0pt}{0.9em}{0.2em}

% Página de cada parte: fondo de color, número grande y una frase de qué contiene.
\newcommand{\parte}[3]{%
  \cleardoublepage\thispagestyle{empty}%
  \addcontentsline{toc}{part}{Parte #1\quad #2}%
  \begin{tikzpicture}[remember picture,overlay]
    \fill[acento] (current page.south west) rectangle (current page.north east);
    \fill[acentoosc] (current page.south west) rectangle ([yshift=7cm]current page.south east);
    \node[anchor=north west,text=white!85!acento,font=\fontsize{120}{120}\selectfont\bfseries]
      at ([xshift=2.3cm,yshift=-3cm]current page.north west) {#1};
    \node[anchor=north west,text=white,font=\fontsize{34}{40}\selectfont\bfseries,text width=15cm,align=left]
      at ([xshift=2.4cm,yshift=-8.4cm]current page.north west) {#2};
    \node[anchor=north west,text=white!90!acento,font=\large,text width=14cm,align=left]
      at ([xshift=2.4cm,yshift=-10.6cm]current page.north west) {#3};
    \node[anchor=south west,text=white!80!acento,font=\small]
      at ([xshift=2.4cm,yshift=2.2cm]current page.south west) {TurbuLENS · Manual de usuario y fundamentos};
  \end{tikzpicture}%
  \clearpage}

% ---------------------------------------------------------------------------------------------- encabezados
\pagestyle{fancy}\fancyhf{}
\fancyhead[L]{\small\color{gris}\textbf{TurbuLENS} 1.5.1}
\fancyhead[R]{\small\color{gris}\leftmark}
\fancyfoot[C]{\small\color{gris}\thepage}
\renewcommand{\headrulewidth}{0.4pt}
\renewcommand{\headrule}{{\color{grisclaro}\hrule width\headwidth height\headrulewidth}}
\renewcommand{\chaptermark}[1]{\markboth{\thechapter.\ #1}{}}
\fancypagestyle{plain}{\fancyhf{}\fancyfoot[C]{\small\color{gris}\thepage}\renewcommand{\headrulewidth}{0pt}}

% ---------------------------------------------------------------------------------------------- índice
\makeatletter
\renewcommand*\l@part[2]{%
  \addpenalty{-\@highpenalty}\vskip 1.1em plus 2pt
  {\setlength\@tempdima{3em}\parindent\z@\rightskip\@pnumwidth\parfillskip-\@pnumwidth
   \leavevmode\large\bfseries\color{acento}#1\hfil\hb@xt@\@pnumwidth{\hss}\par\nobreak}}
\makeatother

% ---------------------------------------------------------------------------------------------- texto de la app
\newcommand{\codigo}[1]{\texttt{#1}}
\newcommand{\ruta}[1]{\texttt{\color{acentoosc}#1}}
\newtcbox{\boton}{on line,colback=white,colframe=gris!45,boxrule=0.5pt,arc=3pt,
  left=3pt,right=3pt,top=1pt,bottom=1pt,boxsep=0pt,fontupper=\small\sffamily}
\newtcbox{\botonp}{on line,colback=acento,colframe=acento,coltext=white,boxrule=0.5pt,arc=3pt,
  left=3pt,right=3pt,top=1pt,bottom=1pt,boxsep=0pt,fontupper=\small\sffamily\bfseries}
\newtcbox{\tecla}{on line,colback=grisclaro,colframe=gris!60,boxrule=0.4pt,arc=2pt,
  left=2.5pt,right=2.5pt,top=0.5pt,bottom=0.5pt,boxsep=0pt,fontupper=\footnotesize\sffamily}
\newcommand{\menu}[1]{\textsf{\color{acentoosc}#1}}
\newcommand{\campo}[1]{\emph{#1}}

% ---------------------------------------------------------------------------------------------- recuadros
\newtcolorbox{nota}[1][Nota]{enhanced,breakable,colback=suave,colframe=acento,boxrule=0pt,leftrule=3pt,arc=0pt,
  left=4mm,right=3mm,top=2mm,bottom=2mm,fonttitle=\bfseries,coltitle=acentoosc,title={#1},
  attach title to upper={\ \ }}
\newtcolorbox{cuidado}[1][Cuidado]{enhanced,breakable,colback=crema,colframe=naranja,boxrule=0pt,leftrule=3pt,arc=0pt,
  left=4mm,right=3mm,top=2mm,bottom=2mm,fonttitle=\bfseries,coltitle=cobre,title={#1},attach title to upper={\ \ }}
\newtcolorbox{consejo}[1][Consejo]{enhanced,breakable,colback=white,colframe=acento!45,boxrule=0.6pt,arc=2mm,
  left=4mm,right=3mm,top=2mm,bottom=2mm,fonttitle=\bfseries,coltitle=acento,title={#1},attach title to upper={\ \ }}
\newtcolorbox{resumen}{enhanced,colback=grisclaro,colframe=grisclaro,arc=2mm,left=5mm,right=5mm,top=3mm,bottom=3mm,
  fontupper=\small,before skip=0.4em,after skip=1em}
\newtcolorbox{pasos}[1][Paso a paso]{enhanced,breakable,colback=white,colframe=acento,boxrule=0.8pt,arc=2mm,
  left=4mm,right=4mm,top=2mm,bottom=2mm,colbacktitle=acento,coltitle=white,fonttitle=\bfseries,title={#1}}

% ---------------------------------------------------------------------------------------------- capturas
\newtcbox{\marcanumcaja}{on line,colback=naranja,colframe=naranja,coltext=white,boxrule=0pt,arc=5.5pt,
  left=0pt,right=0pt,top=0pt,bottom=0pt,boxsep=1.7pt,fontupper=\sffamily\bfseries\scriptsize}
\newcommand{\marcanum}[1]{\marcanumcaja{\makebox[0.62em][c]{\vphantom{8}#1}}}
\newcommand{\marca}[3]{\node[circle,fill=naranja,text=white,draw=white,line width=1pt,inner sep=0pt,
  minimum size=15pt,font=\sffamily\bfseries\footnotesize] at (#1,#2) {#3};}
% \captura[ancho]{archivo}{pie}: captura con borde suave.
\newcommand{\captura}[3][\linewidth]{%
  \begin{figure}[H]\centering
  \tcbox[colback=white,colframe=gris!30,boxrule=0.5pt,arc=1.5mm,left=0pt,right=0pt,top=0pt,bottom=0pt,boxsep=0pt]{%
    \includegraphics[width=#1]{img/#2}}%
  \caption{#3}\label{fig:#2}\end{figure}}
% \capturamarcas[ancho]{archivo}{pie}: lo mismo, con las marcas numeradas de img/<archivo>_marcas.tex.
\newcommand{\capturamarcas}[3][\linewidth]{%
  \begin{figure}[H]\centering
  \tcbox[colback=white,colframe=gris!30,boxrule=0.5pt,arc=1.5mm,left=0pt,right=0pt,top=0pt,bottom=0pt,boxsep=0pt]{%
  \begin{tikzpicture}
    \node[anchor=south west,inner sep=0] (img) at (0,0) {\includegraphics[width=#1]{img/#2}};
    \begin{scope}[x={(img.south east)},y={(img.north west)}]
      \input{img/#2_marcas.tex}
    \end{scope}
  \end{tikzpicture}}%
  \caption{#3}\label{fig:#2}\end{figure}}
% Captura con marcas y su leyenda en un solo bloque, para que nunca queden en páginas distintas:
% \begin{capturam}[ancho]{archivo}{pie}[columnas] \item ... \end{capturam}
\NewDocumentEnvironment{capturam}{O{\linewidth} m m O{2}}
  {\begin{figure}[H]\centering
   \tcbox[colback=white,colframe=gris!30,boxrule=0.5pt,arc=1.5mm,left=0pt,right=0pt,top=0pt,bottom=0pt,boxsep=0pt]{%
   \begin{tikzpicture}
     \node[anchor=south west,inner sep=0] (img) at (0,0) {\includegraphics[width=#1]{img/#2}};
     \begin{scope}[x={(img.south east)},y={(img.north west)}]
       \input{img/#2_marcas.tex}
     \end{scope}
   \end{tikzpicture}}%
   \caption{#3}\label{fig:#2}
   \vspace{-0.3em}\small\RaggedRight
   \begin{multicols}{#4}\begin{enumerate}[label=\marcanum{\arabic*},leftmargin=*,labelsep=0.5em,itemsep=0.25em]}
  {\end{enumerate}\end{multicols}\end{figure}}
% Leyenda suelta: \begin{leyenda} \item ... \end{leyenda}
\newenvironment{leyenda}[1][2]
  {\small\begin{multicols}{#1}\begin{enumerate}[label=\marcanum{\arabic*},leftmargin=*,labelsep=0.5em,itemsep=0.25em]}
  {\end{enumerate}\end{multicols}}

% ---------------------------------------------------------------------------------------------- código
% Generado por scripts/manual_codigo.py: primera línea y archivo de cada fragmento.
\expandafter\def\csname primeratokovinin\endcsname{34}
\expandafter\def\csname ultimatokovinin\endcsname{49}
\expandafter\def\csname archivotokovinin\endcsname{turbulens/core/dimm.py}
\expandafter\def\csname primeraejebase\endcsname{144}
\expandafter\def\csname ultimaejebase\endcsname{166}
\expandafter\def\csname archivoejebase\endcsname{turbulens/core/dimm.py}
\expandafter\def\csname primeraseeingdimm\endcsname{192}
\expandafter\def\csname ultimaseeingdimm\endcsname{211}
\expandafter\def\csname archivoseeingdimm\endcsname{turbulens/core/dimm.py}
\expandafter\def\csname primeraumbral\endcsname{103}
\expandafter\def\csname ultimaumbral\endcsname{117}
\expandafter\def\csname archivoumbral\endcsname{turbulens/core/detect.py}
\expandafter\def\csname primeracentroides\endcsname{130}
\expandafter\def\csname ultimacentroides\endcsname{154}
\expandafter\def\csname archivocentroides\endcsname{turbulens/core/detect.py}
\expandafter\def\csname primeravalido\endcsname{183}
\expandafter\def\csname ultimavalido\endcsname{192}
\expandafter\def\csname archivovalido\endcsname{turbulens/core/detect.py}
\expandafter\def\csname primerasigmaclip\endcsname{46}
\expandafter\def\csname ultimasigmaclip\endcsname{61}
\expandafter\def\csname archivosigmaclip\endcsname{turbulens/core/stats.py}
\expandafter\def\csname primeraescala\endcsname{49}
\expandafter\def\csname ultimaescala\endcsname{53}
\expandafter\def\csname archivoescala\endcsname{turbulens/core/devices.py}
\expandafter\def\csname primeraescalaplaca\endcsname{97}
\expandafter\def\csname ultimaescalaplaca\endcsname{99}
\expandafter\def\csname archivoescalaplaca\endcsname{turbulens/core/devices.py}
\expandafter\def\csname primeraseguimiento\endcsname{327}
\expandafter\def\csname ultimaseguimiento\endcsname{341}
\expandafter\def\csname archivoseguimiento\endcsname{turbulens/core/pipeline.py}
\expandafter\def\csname primeradarkescalado\endcsname{660}
\expandafter\def\csname ultimadarkescalado\endcsname{670}
\expandafter\def\csname archivodarkescalado\endcsname{turbulens/core/calframes.py}
\expandafter\def\csname primeracombinar\endcsname{30}
\expandafter\def\csname ultimacombinar\endcsname{45}
\expandafter\def\csname archivocombinar\endcsname{turbulens/core/calibration.py}
\expandafter\def\csname primeracorregir\endcsname{67}
\expandafter\def\csname ultimacorregir\endcsname{83}
\expandafter\def\csname archivocorregir\endcsname{turbulens/core/calibration.py}
\expandafter\def\csname primeragauss\endcsname{112}
\expandafter\def\csname ultimagauss\endcsname{142}
\expandafter\def\csname archivogauss\endcsname{turbulens/core/starfit.py}
\expandafter\def\csname primerasatformulas\endcsname{85}
\expandafter\def\csname ultimasatformulas\endcsname{99}
\expandafter\def\csname archivosatformulas\endcsname{turbulens/core/satseeing.py}
\expandafter\def\csname primeratendencia\endcsname{102}
\expandafter\def\csname ultimatendencia\endcsname{112}
\expandafter\def\csname archivotendencia\endcsname{turbulens/core/satseeing.py}
\expandafter\def\csname primerasol\endcsname{59}
\expandafter\def\csname ultimasol\endcsname{73}
\expandafter\def\csname archivosol\endcsname{turbulens/core/satseeing.py}
\expandafter\def\csname primeracenitsat\endcsname{51}
\expandafter\def\csname ultimacenitsat\endcsname{56}
\expandafter\def\csname archivocenitsat\endcsname{turbulens/core/satseeing.py}
\expandafter\def\csname primeracarpeta\endcsname{135}
\expandafter\def\csname ultimacarpeta\endcsname{138}
\expandafter\def\csname archivocarpeta\endcsname{turbulens/core/satseeing.py}
\expandafter\def\csname primerarango\endcsname{259}
\expandafter\def\csname ultimarango\endcsname{277}
\expandafter\def\csname archivorango\endcsname{turbulens/core/goes.py}
\expandafter\def\csname primeraproyeccion\endcsname{319}
\expandafter\def\csname ultimaproyeccion\endcsname{332}
\expandafter\def\csname archivoproyeccion\endcsname{turbulens/core/goes.py}

\lstdefinestyle{turbulens}{
  language=Python,
  basicstyle=\ttfamily\footnotesize\color{tinta},
  keywordstyle=\color{acentoosc}\bfseries,
  commentstyle=\color{gris}\itshape,
  stringstyle=\color{cobre},
  numbers=left, numberstyle=\tiny\color{grisnum}, numbersep=8pt,
  columns=fullflexible, keepspaces=true, showstringspaces=false,
  breaklines=true, breakatwhitespace=false, postbreak=\mbox{\textcolor{grisflecha}{$\hookrightarrow$}\space},
  tabsize=4, upquote=true,
  morekeywords={self,True,False,None}
}
\lstset{style=turbulens}
% \fragmento{id}{título}: el fragmento docs/manual/codigo/<id>.py, numerado como en el archivo original.
\newcommand{\fragmento}[2]{%
  \edef\fragprimera{\csname primera#1\endcsname}%
  \edef\fragultima{\csname ultima#1\endcsname}%
  \tcbinputlisting{enhanced,breakable,listing only,listing file={codigo/#1.py},
    listing options={style=turbulens,firstnumber=\fragprimera},
    colback=codigofondo,colframe=grisclaro,boxrule=0.6pt,arc=1.5mm,left=6mm,right=2mm,top=1mm,bottom=1mm,
    colbacktitle=grisclaro,coltitle=tinta,fonttitle=\footnotesize,
    title={\textbf{#2}\hfill\ttfamily\csname archivo#1\endcsname\ ·\ líneas \fragprimera–\fragultima}}}
